% ATTEMPT_STATUS: unresolved
% RESEARCH_GENERATED: 2026-10-01; GPT-6 Astra Ultra; individual mathematical attempt
% TOP_PROBLEM: {"id": 113, "title": "Additive Complements and Cosets", "category_id": 2, "category": {"id": 2, "name": "combinatorics", "display_name": "Combinatorics", "description": "Counting problems, graph theory, discrete structures.", "slug": "combinatorics", "order_index": 2, "created_at": "2026-07-31T15:26:25.671Z"}, "status": "open"}
% FIRST_POSTED: 2026-10-01
% Imported from: unsolvedmath_additional_500.tex; UnsolvedMath ID 113
% !TeX program = lualatex
% Compile twice: lualatex 113.tex
% Self-contained: no external bibliography, images, or \input files.
% Numeric section labels and visible numbers are original UnsolvedMath IDs.
\documentclass[11pt,a4paper]{article}
\usepackage[margin=24mm,headheight=14pt]{geometry}
\usepackage{fontspec}
\setmainfont{Latin Modern Roman}
\setsansfont{Latin Modern Sans}
\setmonofont{Latin Modern Mono}
\usepackage{amsmath,amssymb,mathtools}
\usepackage{unicode-math}
\setmathfont{Latin Modern Math}
\usepackage{microtype}
\usepackage{enumitem,xurl,needspace}
\usepackage[dvipsnames]{xcolor}
\usepackage{hyperref}
\hypersetup{unicode=true,colorlinks=true,linkcolor=MidnightBlue,urlcolor=MidnightBlue,
 pdftitle={UnsolvedMath Problem 113},pdfauthor={Source-annotated compilation},bookmarksnumbered=true}
\usepackage{bookmark,tocloft,titlesec,fancyhdr}
\setlength{\cftsecnumwidth}{4.2em}
\cftsetpnumwidth{2.5em}
\cftsetrmarg{4em}
\titleformat{\section}{\Large\bfseries\raggedright}{\thesection}{0.7em}{}
\pagestyle{fancy}\fancyhf{}
\fancyhead[L]{\small\sffamily UnsolvedMath research attempt}
\fancyhead[R]{\small Original catalogue IDs}
\fancyfoot[C]{\thepage}
\setlength{\parindent}{0pt}
\setlength{\parskip}{5pt plus 1pt}
\setlength{\emergencystretch}{3em}
\setcounter{tocdepth}{1}\setcounter{secnumdepth}{1}
\setlist[itemize]{leftmargin=3em,itemsep=4pt,topsep=4pt,parsep=0pt}
\allowdisplaybreaks
\newcommand{\N}{\mathbb{N}}
\newcommand{\Z}{\mathbb{Z}}
\newcommand{\Q}{\mathbb{Q}}
\newcommand{\R}{\mathbb{R}}
\newcommand{\C}{\mathbb{C}}
\newcommand{\F}{\mathbb{F}}
\newcommand{\A}{\mathbb{A}}
\newcommand{\PP}{\mathbb{P}}
\newcommand{\E}{\mathbb{E}}
\newcommand{\eps}{\varepsilon}
\newcommand{\dd}{\,\mathrm d}
\newcommand{\sref}[2]{\hyperlink{source:#1:#2}{[S#2]}}
\newif\ifshowcatalogue
\showcataloguetrue
\newcommand{\cataloguescope}[1]{\ifshowcatalogue
 \paragraph{Statement recorded in the supplied source.}
 #1\fi}
\begin{document}
% UnsolvedMath ID 113; source catalogue code GREEN-025

\setcounter{section}{112}

\section{Small Additive Complements and Affine Subspaces}\label{113}

{\small\sffamily UnsolvedMath ID 113 · Combinatorics}

\subsection{Definitions and mathematical statement}

\paragraph{Shared notation.}Unless a section says otherwise, $\N=\{1,2,\ldots\}$,
$\Z$ is the ring of integers, $\Q,\R,\C$ are the rational, real, and complex
fields, $[N]=\{1,\ldots,\lfloor N\rfloor\}$, and $\log$ is the natural
logarithm. For real functions of a parameter tending to infinity,
$f\ll g$ and $f=O(g)$ mean $|f|\leq Cg$ eventually for some fixed $C$;
$f\gg g$ reverses this comparison; $f=o(g)$ means $f/g\to0$;
$f\sim g$ means $f/g\to1$; and $f\asymp g$ means both order bounds.
A subscript on $O$ or $\ll$ specifies allowed constant dependence.
The expression $n^{o(1)}$ in an upper bound means growth smaller than
$n^\epsilon$ for each fixed $\epsilon>0$. Local definitions govern any
exception, finite-field convention, or use of the same symbol for another
object. An asymptotic claim is not automatically an effective algorithm.



An additive complement of $A\subseteq\F_2^n$ is a set $B$ such that $A+B=\F_2^n$. The hypothesis is the existence of such a $B$ with $|B|\leq K$. The asserted conclusion is an affine subspace $x+V\subseteq A+A$ with $n-\dim V\leq C(K)$, where $C(K)$ is finite and independent of $n,A,B$. No logarithmic dependence on $K$ is asserted. \sref{113}{1} \sref{113}{2}

\paragraph{Statement recorded in the supplied source.}
Suppose that $A \subset \mathbb{F}_2^n$ has an additive complement of size $K$. Does $2A$ contain a coset of codimension $O_K(1)$? \sref{113}{1}

\paragraph{Residual task reported by the archive.}

The embedded review (2026-08-17) records: Prove or disprove the existence of some codimension bound depending only on K; any valid bound must exceed logarithmic growth along the Hamming-ball family. \sref{113}{1}

\subsection{Short English statement}

If only a bounded number of translates of a set cover the whole binary vector space, must its double sumset contain an affine subspace of bounded codimension?

\subsection{Research attempt}

\paragraph{Provenance and scope.}
This individual attempt was generated by GPT-6 Astra Ultra on 1 October 2026. The method was to derive explicit partial results and test the first proposed extension adversarially. The original statement and sources below are retained. No claim of mathematical novelty or complete solution is made.

\paragraph{Formal target and scope.}
Let $V=\F_2^n$. An additive complement of size at most $K$ means a set $S$ with $|S|\leq K$ and $A+S=V$. The required codimension bound may depend on $K$ but not on $n$. The primary problem list distinguishes this qualitative question from the much stronger logarithmic dependence, and discusses the failure of the latter. Here we prove the cases $K\leq2$, derive an exact colouring reduction, and reconstruct the Hamming-ball obstruction with its covering parameter tracked.

\paragraph{The half-density lemma.}
Every subset $A\subset V$ of cardinality at least $|V|/2$ has $2A$ containing a subspace of codimension at most one. If $|A|>|V|/2$, then for every $t\in V$ the sets $A$ and $A+t$ intersect, giving $t\in2A$; hence $2A=V$. Suppose instead that $|A|=|V|/2$. If $2A=V$ again there is nothing to prove. Otherwise choose $t\notin2A$. Then $A\cap(A+t)=\varnothing$, and equal cardinalities imply $A+t=V\setminus A$.

Let $H=\{h:A+h=A\}$, the translation stabiliser. It is a linear subspace. For any $s$, absence from $2A$ is equivalent to $A+s=V\setminus A=A+t$, which is equivalent to $s+t\in H$. Thus
\[
 V\setminus2A=t+H.
\]
Since $0\in2A$ for nonempty $A$, we have $t\notin H$. Extend a basis of $H$ by $t$ and choose a linear functional $\ell$ vanishing on $H$ with $\ell(t)=1$. The hyperplane $\ker\ell$ avoids $t+H$ and is contained in $2A$. This proves the lemma, including the affine-coset issue.

If $|S|\leq2$, the covering inequality $|V|=|A+S|\leq2|A|$ gives the needed half density. For $K=1$, $A$ is a translate of the whole space and hence equals $V$; for $K=2$, codimension at most one follows. The bound is sharp for a hyperplane $A$ with a two-element complement meeting both cosets.

\paragraph{Reduction from partitions, with translations cancelled.}
Enumerate $S=\{s_1,\ldots,s_K\}$. Assign each $x\in V$ to the first index $i$ for which $x\in A+s_i$, obtaining a genuine partition $V=C_1\sqcup\cdots\sqcup C_K$ with $C_i\subset A+s_i$. In characteristic two,
\[
 C_i+C_i\subset(A+s_i)+(A+s_i)=2A.
\]
Therefore a theorem giving a bounded-codimension coset in the double sumset of some colour class in every $K$-colour partition would imply the present covering theorem with the same bound. This proves the implication exactly. Choosing the largest $C_i$ and using density alone does not prove it: density at least $1/K$ does not itself force bounded codimension in a two-fold sumset.

\paragraph{Why logarithmic dependence fails in this method's natural examples.}
Take $A=B_r$, the Hamming ball with $r=n/2-\lfloor\sqrt n\rfloor$ for even $n$. Its density is bounded below by an absolute positive constant $a$; this follows by summing the order-$2^n/\sqrt n$ central binomial coefficients in a window of width order $\sqrt n$ lying below $r$. Its double sumset is exactly $B_{2r}$, because a support of size at most $2r$ splits into two supports of size at most $r$.

Every affine subspace of dimension $d$ contains a point of weight at least $d$: choose $d$ pivot coordinate columns in a basis matrix and prescribe ones on those coordinates. Consequently every affine subspace in $2A$ has dimension at most $2r$, so its codimension is at least $2\lfloor\sqrt n\rfloor$.

Nevertheless $A$ has an additive complement of size $O(n)$. Choose $k$ independent uniform translations. For each fixed point the probability of remaining uncovered is $(1-\alpha)^k\leq e^{-ak}$. A union bound over $2^n$ points is less than one if $k>(n\log2+1)/a$. Some choice therefore covers the space. Repeated translations can be removed, so an actual set of at most $k$ translations works. Since $\sqrt n$ grows faster than $\log k$, no universal $O(\log K)$ conclusion can hold for these examples.

\paragraph{Verification and exact remaining gap.}
The random-cover argument proves existence, not an efficient covering algorithm, and it keeps the ambient dimension in the number of translations. Thus it does not refute a bound depending arbitrarily on a fixed $K$. The stabiliser argument uses exact half density and cannot be applied unchanged to density $1/3$. The work proves the full $K\leq2$ cases, the partition implication, and a precise obstruction to logarithmic strengthening. The unresolved step is a dimension-independent codimension bound for every fixed $K\geq3$.

\subsection{Sources}

{\small

\begin{itemize}

\item[\hypertarget{source:113:1}{S1}] ULAM.AI, \emph{UnsolvedMath}, supplied \texttt{problems.json}, record 113 (GREEN-025), statement and background. The embedded literature-review date is 2026-08-17; that is the archive’s review date, not a fresh certification. Dataset landing page: \url{https://huggingface.co/datasets/ulamai/UnsolvedMath}.

\item[\hypertarget{source:113:2}{S2}] Ben Green, 100 Open Problems, Problem 52 and 2025 update, . \url{https://people.maths.ox.ac.uk/greenbj/papers/open-problems.pdf}. The author-hosted document was inspected on 27 September 2026; the archive’s local code is not a problem-number locator in this paper.

\end{itemize}

}

\label{end:113}
\end{document}
